\section{Discussion and limitations}
\paragraph{Composition, not risk.} The outcome is conditional on an incident having been reported as one of four crime types. A falling robbery odds ratio can come from fewer robberies, more of the other crimes or changes in reporting and classification; the data cannot separate these, and our composition check shows that the pooled post-2019 ratio depends heavily on the comparison class. None of the estimates is a probability of being robbed or an effect of time on crime.

\paragraph{Premises are nearly definitional.} The largest predictive lift comes from premises and location type, which partly define the crime types (a break and enter occurs at a building). That lift is therefore not a discovery about robbery; it is why we report the no-premises ablation. Without those variables, even the best model has only modest predictive value.

\paragraph{Census features are coarse.} The census block is a single profile that is constant within a neighbourhood and over the ten years, so it cannot track change within a neighbourhood and can proxy for neighbourhood identity. Grouped cross-validation prevents that from leaking across folds, and the block's small gain should be read with this in mind. Across neighbourhoods, 2020 robbery counts relate only weakly to median income (Spearman $\rho =$ \res{predict.census_scatter:summary.spearman_income.rho}) and not to the job ratio ($\rho =$ \res{predict.census_scatter:summary.spearman_ratio.rho}); an interactive three-dimensional view is on the project site.

\paragraph{Data quality and shift.} Police-recorded incidents reflect reporting and recording practices, which differ by crime type and neighbourhood \citep{skogan1975dark}, and the published coordinates are offset by the publisher, so geography is approximate. The robbery share moved from \resG{prevalence_by_year.2014.share} to \resG{prevalence_by_year.2023.share}, a prior shift that lowers precision at a fixed threshold in the temporal holdout and means thresholds tuned in one period do not transfer unchanged. Because F1 thresholds are tuned on inner out-of-fold predictions from models that saw less data (and XGBoost was early-stopped on the scored fold), F1, precision and recall are slightly optimistic.

\paragraph{Significance at scale and scope.} With \resG{n} rows, intervals are narrow, and a statistically clear term need not be large or useful. We also limit the study to four crime types and do not model assault, bicycle theft, homicide or shootings. Predictive scores answer only which type a reported incident is, not where or when to expect incidents.
